% !PRINTED 102
\weberpage{102}

Wir sagen auch kurz, $\alpha$ ist eine Wurzel von $f(x)$. Nach §. 4 lässt sich dann $f(x)$ durch $x-\alpha$ ohne Rest theilen, so dass man
\begin{equation}
 f(x)=(x-\alpha)f_1(x)
\tag{2}
\end{equation}
setzen kann, worin $f_1(x)$ nur vom $(n-1)$ten Grad ist, und denselben ersten Coëfficienten $a_0$ hat, wie $f(x)$, also
\[
 f_1(x)=a_0x^{n-1}+a'_1x^{n-2}+\cdots .
\]
Die Coëfficienten von $f_1(x)$ sind in §. 4 (6) angegeben.

Jede Wurzel von $f_1(x)$ ist also zugleich Wurzel von $f(x)$; und umgekehrt ist jede Wurzel von $f(x)$ entweder $=\alpha$ oder eine Wurzel von $f_1(x)$. Wenn eine Wurzel $\alpha$ von $f(x)$ bekannt ist, so ist die Aufgabe, die übrigen zu finden, auf die Lösung einer Gleichung $(n-1)$ten Grades zurückgeführt.

Das Ziel der Betrachtungen dieses Abschnittes besteht in dem Nachweis, dass jede Function $f(x)$ vom $n$ten Grad wenigstens eine Wurzel hat. Dieser Satz heisst der \gesperrt{Fundamentalsatz der Algebra}. Zunächst ziehen wir aus (2) die wichtige Folgerung:

\medskip
\noindent I. \gesperrt{Eine Gleichung $n$ten Grades kann nicht mehr als $n$ Wurzeln haben.}
\medskip

Denn hätte $f(x)$ mehr als $n$ Wurzeln, so hätte $f_1(x)$, was nur vom $(n-1)$ten Grade ist, mehr als $n-1$ Wurzeln. Eine Gleichung ersten Grades hat aber nicht mehr als eine Wurzel, woraus die Richtigkeit unseres Satzes durch vollständige Induction folgt.

Man giebt ihm bisweilen auch den Ausdruck:

\medskip
\noindent II. \gesperrt{Wenn eine Function $n$ten Grades $f(x)$ mehr als $n$ Wurzeln hat, so müssen alle ihre Coëfficienten Null sein.}
\medskip

Ist $\beta$ eine Wurzel von $f_1(x)$, so lässt sich ebenso setzen
\[
 f_1(x)=(x-\beta)f_2(x),\quad f(x)=(x-\alpha)(x-\beta)f_2(x),
\]
worin $f_2(x)=a_0x^{n-2}+\cdots$ nur vom $(n-2)$ten Grade ist. Nimmt man also an, dass jede der Functionen, die man durch diese Division erhält, $f_1(x), f_2(x)\ldots$ wenigstens eine Wurzel habe, so erhält man schliesslich
\begin{equation}
 f(x)=a_0(x-\alpha)(x-\beta)\cdots(x-\nu),
\tag{3}
\end{equation}
und wir können also den Satz, den wir vorhin als das Ziel unserer Betrachtungen bezeichnet haben, auch so aussprechen: Es soll bewiesen werden:

\medskip
